% Corrected French for printed p.167.
% The printed coho-/mologiques division is retained; the missing inline support Phi is corrected.
\SourcePage{167}

\LayerTarget{tohoku:construction:cohomology-support-comparison:p2}
\noindent mologiques se réduisant en dimension 0 à celui envisagé plus haut) ; ou encore,
partant d'un faisceau $F$ sur $X$ :

\LayerTarget{tohoku:eq:canonical-support-comparison-FZ}
\[
H^p_{\Phi_Z}(Z,F\vert Z)\to H^p_\Phi(X,F_Z).
\]

\LayerTarget{tohoku:theorem:3.5.1}
\indent\textsc{Théorème 3.5.1.} \textit{Les homomorphismes précédents
sont des isomorphismes dans chacun des deux cas suivants :}
\begin{enumerate}
\item \textit{$Z$ est fermé.}
\item \textit{$\Phi$ est paracompactifiante, $Z$ est ouvert.}
\end{enumerate}

\LayerTarget{tohoku:proof:3.5.1}
\indent\textsc{Démonstration.} Il suffit dans chacun de ces deux cas, de
vérifier que les foncteurs $H^p_\Phi(X,G^X)$ sur $\mathbf C^Z$ sont effaçables pour
$p\geqq1$. Si $Z$ est fermé, cela résulte du fait que si $G$ est injectif,
$G^X$ est injectif (prop. 3.1.4). Si $Z$ est ouvert, et si $\Phi$ est
paracompactifiante, il en est de même de $\Phi_Z$, donc tout $G\in \mathbf C^Z$ se
plonge dans un faisceau $\Phi_Z$-mou (cf. 3.3), il suffit alors de noter que si
$G$ est $\Phi_Z$-mou, alors $G^X$ est $\Phi$-mou, (fait dont la vérification
est immédiate), et d'appliquer le corollaire à prop. 3.3.2.

Revenons alors à la suite exacte
$0\to F_U\to F\to F_Y\to0$
relative à un sous-espace fermé $Y$ et son complémentaire ouvert $U$, elle
donne naissance à une suite exacte de cohomologie, qu'on peut écrire, grâce au
théorème précédent :

\LayerTarget{tohoku:eq:closed-open-support-long-exact-sequence}
\[
\mathinner{\cdotp\cdotp\cdotp\cdotp} H^p_\Phi(X,F_U)\to H^p_\Phi(X,F)
\to H^p_{\Phi_Y}(Y,F)\to
H^{p+1}_\Phi(X,F_U)\mathinner{\cdotp\cdotp\cdotp\cdotp}
\]
(où pour simplifier, on a écrit $F$ au lieu de $F\vert Y$ dans le troisième
terme).
\LayerTarget{tohoku:eq:open-support-identification}
Si $\Phi$ est paracompactifiante, on peut de plus remplacer les termes de la
forme $H^p_\Phi(X,F_U)$ par $H^p_{\Phi_U}(U,F)$, et on obtient alors la suite
exacte bien connue de [4, XVII].

\LayerTarget{tohoku:section:3.6}
\indent\textbf{3.6. Sur la dimension cohomologique de certains espaces}%
\LayerTarget{tohoku:footnote:8}\textsuperscript{8)}.

\LayerTarget{tohoku:proposition:3.6.1}
\indent\textsc{Proposition 3.6.1.} \textit{Soient $X$ un espace
topologique, $(T^p)$ un foncteur cohomologique covariant, défini sur la
catégorie $\mathbf C^X$ des faisceaux de groupes abéliens sur $X$, à valeurs dans une
catégorie $\mathbf C'$. On suppose que $\mathbf C'$ satisfait la condition \textup{AB 4)} \textup{(}cf. \textup{1.5)} qui
implique que l'on peut former dans $\mathbf C'$ des limites inductives \textup{(}prop. \textup{1.8)}, et
on suppose que les $T^p$ permutent aux limites inductives. Soit $F\in \mathbf C^X$,
alors $T^p(F)$ appartient à toute sous-catégorie épaisse \textup{(}cf. \textup{1.11)} $\mathbf C''$ de
$\mathbf C'$ stable par sommes directes infinies, qui contient tous les objets de la
forme $T^i(\mathbf Z_U)$, où $U$ est un ouvert arbitraire de $X$ et où $i$ est égal à
$p,p+1$ ou $p+2$.}

\LayerTarget{tohoku:proof:3.6.1:p1}
(La signification de $\mathbf Z_U$ est la même qu'au numéro précédent).
\LayerTarget{tohoku:construction:generator-family}
Considérons une famille $(f_i)_{i\in I}$ de sections de $F$ sur des ouverts
$U_i$ ; chaque $f_i$ définit un homomorphisme de $\mathbf Z_{U_i}$ dans $F$, donc
$(f_i)$ définit un homomorphisme de la
\LayerTarget{tohoku:eq:generator-epimorphism}
somme directe $\displaystyle\bigoplus_i \mathbf Z_{U_i}$ dans $F$.
\LayerTarget{tohoku:def:system-of-generators}
Nous dirons que $(f_i)$ est un \textit{système de générateurs} de $F$ si
l'homomorphisme précédent est un épimorphisme. Il est trivial qu'il existe
toujours, pour $F$ donné, une famille de générateurs, d'où résulte aussitôt que
$F$ est limite inductive d'une famille filtrante croissante

\TohokuContinuation{168}
